\begin{frame}{Proving IFF}
\begin{center}
\vspace{1.5cm}
{\Huge\textbf{Proving IFF}}
\end{center}
\end{frame}

% --------------------------------------------------
% Question 1
% --------------------------------------------------

\begin{frame}{Question 1}
A software system claims that an integer $n$ is even if and only if
$n^2$ is even.

Prove the claim.
\end{frame}

\begin{frame}{Question 1 -- Answer}
To prove:

$$
n\text{ is even}\iff n^2\text{ is even}
$$

we must prove both directions.

First:

$$
n\text{ even}\Rightarrow n^2\text{ even}
$$

\end{frame}

\begin{frame}{Question 1 -- Answer}
Let:

$$
n=2k
$$

Then:

$$
n^2=(2k)^2=4k^2
$$

$$
n^2=2(2k^2)
$$

Therefore, $n^2$ is even.
\end{frame}

\begin{frame}{Question 1 -- Answer}
For the reverse direction, assume:

$$
n^2\text{ is even}
$$

Suppose $n$ were odd. Then:

$$
n=2k+1
$$

and:

$$
n^2=4k^2+4k+1
$$

which is odd.

This is a contradiction.

Therefore, $n$ is even.
\end{frame}

\begin{frame}{Question 1 -- Answer}
Both directions have been proved.

Hence:

$$
\boxed{n\text{ is even}\iff n^2\text{ is even}}
$$

\end{frame}

% --------------------------------------------------
% Question 2
% --------------------------------------------------

\begin{frame}{Question 2}
A database stores an integer value $n$. A validation rule states that $n$
is divisible by $6$ if and only if it is divisible by both $2$ and $3$.

Prove the rule.
\end{frame}

\begin{frame}{Question 2 -- Answer}
We prove both directions.

First:

$$
6\mid n\Rightarrow(2\mid n\land3\mid n)
$$

If:

$$
n=6k
$$

then:

$$
n=2(3k)=3(2k)
$$

Therefore, $2\mid n$ and $3\mid n$.
\end{frame}

\begin{frame}{Question 2 -- Answer}
For the reverse direction, assume:

$$
2\mid n
$$

and:

$$
3\mid n
$$

Since $2$ and $3$ are relatively prime, their product divides $n$.

Therefore:

$$
6\mid n
$$

\end{frame}

\begin{frame}{Question 2 -- Answer}
Thus both directions hold:

$$
\boxed{6\mid n\iff(2\mid n\land3\mid n)}
$$

\end{frame}

% --------------------------------------------------
% Question 3
% --------------------------------------------------

\begin{frame}{Question 3}
A programming system uses an integer $n$. Prove that $n$ is odd if and only
if $n+1$ is even.
\end{frame}

\begin{frame}{Question 3 -- Answer}
First direction:

$$
n\text{ odd}\Rightarrow n+1\text{ even}
$$

Let:

$$
n=2k+1
$$

Then:

$$
n+1=2k+2=2(k+1)
$$

Therefore, $n+1$ is even.
\end{frame}

\begin{frame}{Question 3 -- Answer}
Reverse direction:

$$
n+1\text{ even}\Rightarrow n\text{ odd}
$$

Let:

$$
n+1=2k
$$

Then:

$$
n=2k-1
$$

$$
n=2(k-1)+1
$$

Therefore, $n$ is odd.
\end{frame}

\begin{frame}{Question 3 -- Answer}
Both directions hold.

Hence:

$$
\boxed{n\text{ is odd}\iff n+1\text{ is even}}
$$

\end{frame}

% --------------------------------------------------
% Question 4
% --------------------------------------------------

\begin{frame}{Question 4}
A network system uses two integer parameters $a$ and $b$.

Prove that:

$$
a+b\text{ is even}
$$

if and only if $a$ and $b$ have the same parity.
\end{frame}

\begin{frame}{Question 4 -- Answer}
First assume $a$ and $b$ have the same parity.

If both are even:

$$
a=2m,\qquad b=2n
$$

Then:

$$
a+b=2(m+n)
$$

\end{frame}

\begin{frame}{Question 4 -- Answer}
If both are odd:

$$
a=2m+1,\qquad b=2n+1
$$

Then:

$$
a+b=2(m+n+1)
$$

Thus, in either case, $a+b$ is even.
\end{frame}

\begin{frame}{Question 4 -- Answer}
Now assume:

$$
a+b\text{ is even}
$$

If $a$ and $b$ had different parity, one would be even and the other odd.

Their sum would then be odd.

This contradicts the assumption.

Therefore, $a$ and $b$ have the same parity.

Hence:

$$
\boxed{a+b\text{ is even}\iff
a\text{ and }b\text{ have the same parity}}
$$

\end{frame}

% --------------------------------------------------
% Question 5
% --------------------------------------------------

\begin{frame}{Question 5}
A data structure contains an integer number of elements $n$.

Show that $n(n-1)$ is even if and only if at least one of $n$ and $n-1$
is even.
\end{frame}

\begin{frame}{Question 5 -- Answer}
First, suppose at least one of $n$ and $n-1$ is even.

Since:

$$
n(n-1)
$$

is a product containing an even factor, the product is even.
\end{frame}

\begin{frame}{Question 5 -- Answer}
For the reverse direction, suppose:

$$
n(n-1)\text{ is even}
$$

Two consecutive integers cannot both be odd.

Therefore, among $n$ and $n-1$, at least one must be even.

Thus:

$$
\boxed{n(n-1)\text{ is even}
\iff
n\text{ or }n-1\text{ is even}}
$$

\end{frame}
